% ============================================================================
%  RELATED WORK -- drop-in section for the ASIA CCS / ACM CCS paper.
%
%  Assumptions made while writing this (check before using):
%    - Placed after the Introduction, before Threat Model / Preliminaries.
%    - "our attack" is referenced only as \ourmethod{} and by a forward
%      \ref{sec:method}; the macro is defined below so you can bind it to
%      whichever name (VLI*, ZLI*, ZLI) the rest of the paper has settled on
%      without editing this file again.
%    - No .bib file exists yet in the repo; a plain bibitem list matching
%      every \cite key used here is given at the end. Merge it into your
%      bibliography (or convert to BibTeX) before compiling.
%    - iLRG is cited under its official name (the paper's own GitHub repo is
%      BUAA-CST/iLRG); "iRLG", used in this codebase and in RLU's own related
%      work section, is a naming carry-over this paper does not repeat.
% ============================================================================

\newcommand{\ourmethod}{ZLI}

\section{Related Work}
\label{sec:related}

\paragraph{Gradient inversion attacks.}
Federated learning lets clients contribute to a shared model by exchanging
gradients or weight updates rather than raw data~\cite{mcmahan2017fedavg}.
DLG~\cite{zhu2019dlg} showed this exchange is not safe: matching a dummy
input and a relaxed, continuous label distribution against the observed
gradient recovers both, though the joint search degrades as batch size and
resolution grow. Follow-up work improved the image-reconstruction half of
this search with cosine-similarity objectives and total-variation
regularization, reaching ImageNet resolution~\cite{geng2022zlg}, but every
such method still pays for the label half of the search, either by
optimizing over it jointly or by avoiding batches with repeated labels.
This motivated recovering the label distribution first, in closed form,
and using it to constrain reconstruction afterward --- reported to improve
reconstructed image quality in its own right~\cite{ma2023ilrg,chen2024rlu}.
We are concerned entirely with that first stage: recovering the per-class
\emph{instance count} of a batch from a shared update, independent of
whichever image-reconstruction attack follows it.

\paragraph{Label-distribution inference attacks.}
A series of closed-form estimators recover per-class counts directly from
the sign and magnitude of the final layer's gradient, without an
optimization loop over inputs. iDLG~\cite{zhao2020idlg} first showed that a
single sample's bias gradient is negative only at its ground-truth class,
but only at batch size one; GradInversion~\cite{yin2021gradinversion}
extended the same sign rule to full batches using gradient magnitude, at
the cost of assuming no two samples share a label. Two methods estimate
their per-class statistics from the attacked update alone: LLG~\cite{wainakh2022llg}
derives a per-class impact and offset analytically from the single observed
gradient and peels instances off greedily; iLRG~\cite{ma2023ilrg} reconstructs
a per-class embedding from the final layer's gradients and feeds it back
through the \emph{attacked model's own} classifier to approximate
post-softmax probabilities, solving the resulting linear system via the
Moore--Penrose pseudo-inverse. The method of Geng et al.~\cite{geng2022zlg},
referred to as ZLG in follow-up work though not named as such in the
original paper, is likewise self-contained in its original form: it
estimates the same kind of per-class confidence and mean embedding by
passing \emph{dummy}, synthetically generated images through the current
model rather than any real dataset.

Three variants strengthen these estimates with a small auxiliary dataset
the server is assumed to hold, drawn from a distribution similar to the
clients' data. Wainakh et al.\ introduce this variant themselves: LLG+~\cite{wainakh2022llg}
re-estimates LLG's impact and offset by running the current model on real
auxiliary samples in place of dummy ones. ZLG+ is not Geng et al.'s own
design; it was introduced by Chen \& Vikalo~\cite{chen2024rlu} as a fairness
baseline for RLU, substituting the same real auxiliary set for ZLG's
dummy-data estimation step so that every baseline in their comparison draws
on the same side information. RLU~\cite{chen2024rlu} itself poses recovery
as a least-squares problem whose per-class prior is obtained by Monte Carlo
evaluation of the auxiliary set, and is, to our knowledge, the first in
this line to remain effective across multiple local epochs, optimizers, and
non-IID partitions. Because auxiliary dependence enters ZLG+ as a
benchmarking choice rather than a property of the original attack, we
evaluate ZLG under both its original (dummy-data) and auxiliary-data
(ZLG+) forms throughout this paper. RLG~\cite{dang2021rlg} recovers only the
\emph{set} of labels present in a batch, across a broader range of
architectures and modalities including speech; because it targets class
existence rather than per-class counts, it has no counterpart to the
instance-level recovery rate we use throughout this paper, and we exclude
it from our count-level comparison.

All three auxiliary-dependent baselines above are reported to degrade as
the attacked model converges, which their own authors attribute to
shrinking gradient signal. Section~\ref{sec:method} argues this account is
incomplete: each of these attacks recovers counts by comparing the
observed update against a center derived, in whole or in part, from
auxiliary statistics, and that center does not shrink along with the
update as training progresses. The two quantities are tied together only
by the gradient's own magnitude, so as that magnitude falls the recovered
distribution is increasingly determined by the auxiliary center rather
than by the client's data. This scale mismatch, not merely the shrinking
signal itself, is what motivates \ourmethod{}, an attack whose recovery
rule is invariant to the scale of the update and therefore does not
require an auxiliary dataset to supply that center at all.

% ----------------------------------------------------------------------------
% Bibliography entries for the \cite keys used above.
% No .bib exists in the repo yet -- merge these into one, or convert to
% \bibitem form inside your existing \begin{thebibliography}.
% ----------------------------------------------------------------------------
% @inproceedings{mcmahan2017fedavg,
%   author = {McMahan, H. Brendan and Moore, Eider and Ramage, Daniel and Hampson, Seth and Ag{\"u}era y Arcas, Blaise},
%   title  = {Communication-Efficient Learning of Deep Networks from Decentralized Data},
%   booktitle = {AISTATS},
%   year   = {2017}
% }
% @inproceedings{zhu2019dlg,
%   author = {Zhu, Ligeng and Liu, Zhijian and Han, Song},
%   title  = {Deep Leakage from Gradients},
%   booktitle = {NeurIPS},
%   year   = {2019}
% }
% @article{zhao2020idlg,
%   author  = {Zhao, Bo and Mopuri, Konda Reddy and Bilen, Hakan},
%   title   = {iDLG: Improved Deep Leakage from Gradients},
%   journal = {arXiv:2001.02610},
%   year    = {2020}
% }
% @inproceedings{yin2021gradinversion,
%   author    = {Yin, Hongxu and Mallya, Arun and Vahdat, Arash and Alvarez, Jose M. and Kautz, Jan and Molchanov, Pavlo},
%   title     = {See through Gradients: Image Batch Recovery via GradInversion},
%   booktitle = {CVPR},
%   year      = {2021}
% }
% @inproceedings{wainakh2022llg,
%   author    = {Wainakh, Aidmar and Ventola, Fabrizio and M{\"u}{\ss}ig, Till and Keim, Jens and Garcia Cordero, Carlos and Zimmer, Ephraim and Grube, Tim and Kersting, Kristian and M{\"u}hlh{\"a}user, Max},
%   title     = {User-Level Label Leakage from Gradients in Federated Learning},
%   booktitle = {PETS},
%   year      = {2022},
%   note      = {arXiv:2105.09369}
% }
% @inproceedings{geng2022zlg,
%   author    = {Geng, Jiahui and Mou, Yongli and Li, Feifei and Li, Qing and Beyan, Oya and Decker, Stefan and Rong, Chunming},
%   title     = {Towards General Deep Leakage in Federated Learning},
%   booktitle = {arXiv:2110.09074},
%   year      = {2022}
% }
% @inproceedings{ma2023ilrg,
%   author    = {Ma, Kailang and Sun, Yu and Cui, Jian and Li, Dawei and Guan, Zhenyu and Liu, Jianwei},
%   title     = {Instance-wise Batch Label Restoration via Gradients in Federated Learning},
%   booktitle = {ICLR},
%   year      = {2023}
% }
% @inproceedings{chen2024rlu,
%   author    = {Chen, Huancheng and Vikalo, Haris},
%   title     = {Recovering Labels from Local Updates in Federated Learning},
%   booktitle = {ICML},
%   year      = {2024},
%   note      = {arXiv:2405.00955}
% }
% @inproceedings{dang2021rlg,
%   author    = {Dang, Trung and Thakkar, Om and Ramaswamy, Swaroop and Mathews, Rajiv and Chin, Peter and Beaufays, Fran{\c c}oise},
%   title     = {Revealing and Protecting Labels in Distributed Training},
%   booktitle = {NeurIPS},
%   year      = {2021},
%   note      = {arXiv:2111.00556}
% }
